\section{Legal logic: the fixed-premise inference boundaries}
\label{annex:p11-rat113-fixed-premises}

\begin{annexitem}{annex:p11.rat113-detachment}{thm:p11-rat113-detachment}{Conditional detachment without an automatic converse}
\begin{proof}
In any classical first-order structure satisfying
\(\forall x(T(x)\to R(x))\), the conditional
\(T(a)\to R(a)\) holds at the interpretation of the constant \(a\).
If \(T(a)\) also holds, its consequent \(R(a)\) holds by the truth
condition for material implication. Thus every model of both premises is
a model of the conclusion.

To show that the fact alone is insufficient, use a one-object domain
\(\{d\}\), interpret \(a\) by \(d\), and set \(T(d)\) true and
\(R(d)\) false. This satisfies \(T(a)\) and falsifies \(R(a)\).
It does not satisfy the rule, precisely the missing normative premise.

For the nonconverse boundary, use the same domain and constant but set
\(T(d)\) false and \(R(d)\) true. The universal conditional is true
because its antecedent is false, and \(\neg T(a)\) is true, while
\(\neg R(a)\) is false. A sufficient condition is not a necessary one.

Finally, if every model of a premise set \(\Gamma\) satisfies a
conclusion and \(\Gamma\subseteq\Delta\), every model of
\(\Delta\) is also a model of \(\Gamma\). It therefore satisfies
the same conclusion. This is monotonicity of the fixed classical consequence
relation, including the vacuous case in which \(\Delta\) has no model.
It says nothing about whether newly discovered information warrants
retaining the old interpreted legal premise. Revising that premise or
using a separately specified defeasible relation changes the inference
problem; it does not refute the conditional truth condition.
\end{proof}
\end{annexitem}

\begin{annexitem}{annex:p11.rat113-extension-boundaries}{thm:p11-rat113-extension-boundaries}{Countermodels and the missing transfer/exclusivity premises}
\begin{proof}
Fix \(n\geq1\). Take a one-object domain \(\{d\}\) and interpret
\(a\) by \(d\). Set \(T_1(d)\) false, \(T'_1(d)\) true,
every \(T_i(d)\) for \(2\leq i\leq n\) true, and \(R(d)\)
false. For \(n=1\), the conjunction of the remaining conditions is
true by the declared convention. Thus \(T(d)\) is false and
\(T'(d)\) is true. The original universal conditional is true, the
comparison-case premise \(T'(a)\) is true, but \(R(a)\) is false.
This proves failure of the unaugmented analogical inference in every
stated case, including the one-condition boundary.

If \(\forall x(T'(x)\to T(x))\) is added, its instance at \(a\)
and \(T'(a)\) yield \(T(a)\); the original rule then yields
\(R(a)\). If instead \(\forall x(T'(x)\to R(x))\) is supplied,
its instance and \(T'(a)\) yield \(R(a)\) directly. These are
different possible justifications: inclusion under the old conditions
versus extension of the consequence to different conditions.

For contrary reasoning, the one-object model with \(T(d)\) false and
\(R(d)\) true again satisfies the original rule and \(\neg T(a)\)
but falsifies \(\neg R(a)\). Now add
\(\forall x(R(x)\to T(x))\). In any model satisfying this added
premise and \(\neg T(a)\), \(R(a)\) cannot be true, for then the
instance \(R(a)\to T(a)\) would force \(T(a)\), contradicting the
negative factual premise. Hence \(\neg R(a)\) follows. The reverse
implication supplies the missing necessity/exclusivity condition.

No countermodel here is asserted to describe a real jurisdiction; it tests
what follows from the specified formal premises. Nor does the validity
after augmentation prove that a legislature, precedent or admissible
interpretative method authorizes the added premise. That is the distinct
legal-methodological question the formal reconstruction exposes.
\end{proof}
\end{annexitem}
